\section{A non-GRS pencil and its quotient}
\label{sec:pencil}

We next give a substantial six-point application.  The family below contains
both conic and nonconic members and retains nontrivial projective variation
away from the conic divisor.  It is useful because its projective equivalence
problem has an explicit one-dimensional quotient, so the geometric branch of
the all-length theorem can be seen in a complete family rather than at an
isolated code.

Over a finite field of odd order, put
\[
H_t=\bigl((0,1,1-t),(0,1,t-1),(1,1-t,0),(1,t-1,0),
          (1,0,-t),(1,0,t)\bigr),
\]
and define
\[
 A(t)=-4t(t-1)^2,\qquad
 B(t)=(t^2-t+1)(t^2-3t+1),\qquad
 z(t)=\bigl(B(t)/A(t)\bigr)^2.
\]
The form of \(z\) is dictated by the residual parameter symmetries.  The
intermediate coordinate
\[
 y(t)=\frac{(t-1)^2}{t}
\]
is generically two-to-one in \(t\), and projective equivalence identifies
\(y\) with \(\pm y^{\pm1}\).  Since
\(z=(y-y^{-1})^2/16\), the map \(t\mapsto z\) is generically eight-to-one.

The admissibility conditions have separate meanings.  The factors
\(t(t-1)(t^2-t+1)(t^2-3t+1)\) exclude arc degenerations and singularities of
the displayed quotient coordinates.  The factor
\[
 G(t)=t^4-4t^3+7t^2-4t+1
\]
cuts out the conic, equivalently GRS, locus.  Define the regular non-GRS set
\begin{equation}\label{eq:regular-pencil-locus}
 \mathcal U=
 \{t:t(t-1)(t^2-t+1)(t^2-3t+1)G(t)\ne0\}.
\end{equation}
Factoring the twenty \(3\)-by-\(3\) column minors gives the first four
factors, while the determinant for a conic through the six columns is a unit
times \(G\).  For \(t\in\mathcal U\), write
\(C_t=\ker H_t\) and \(\Psi_t=\Psi_{C_t}\).

Here \(\sim_{\mathrm{proj}}\) means projective equivalence of the six-point
configurations after independent column rescaling and coordinate permutation;
\(\sim_{\mathrm{mon}}\) is monomial code equivalence; and
\(\sim_{\LC}\), \(\sim_{\LU}\) allow tensor-factor permutations followed by
products of one-qudit Clifford or unitary operators, respectively.

\begin{theorem}[Projective and Clifford quotient]
\label{thm:lc-pencil}
For \(t,u\in\mathcal U\), the following are equivalent:
\[
H_t\sim_{\mathrm{proj}}H_u,\qquad
C_t\sim_{\mathrm{mon}}C_u,\qquad z(t)=z(u),
\]
allowing arbitrary relabelings of the six coordinates.  If \(q\) is prime, these are
also equivalent to \(\Psi_t\sim_{\LC}\Psi_u\).
\end{theorem}

\begin{corollary}[Quantum classification on the regular locus]
\label{cor:lu-lc-pencil}
If \(q\) is an odd prime and \(t,u\in\mathcal U\), then
\[
 H_t\sim_{\mathrm{proj}}H_u
 \Longleftrightarrow C_t\sim_{\mathrm{mon}}C_u
 \Longleftrightarrow \Psi_t\sim_{\LC}\Psi_u
 \Longleftrightarrow \Psi_t\sim_{\LU}\Psi_u
 \Longleftrightarrow z(t)=z(u),
\]
where the code and point relabeling is the corresponding tensor-factor
permutation for the quantum states.
\end{corollary}

\AMEFigureSix

\begin{proof}[Proof of Theorem~\ref{thm:lc-pencil}]
The proof has two independent recovery steps.  First, equality of \(z\) is
equivalent to the residual orbit relation
\(y(u)\in\{\pm y(t)^{\pm1}\}\).  Each of these four relations is realized by
an explicit projectivity and coordinate permutation.  For the converse,
write \([ijk]\) for the determinant of columns \(i,j,k\).  The ten signed
products \([ijk][\ell mn]\) over complementary triples have multiplicity
pattern
\[
 \{+A,+A,+A,-A,-A,-A,+B,+B,-B,-B\},
\]
which recovers \((B/A)^2=z\), including the collision \(z=1\).  This proves
the projective classification and, through the MDS dictionary of
Section~\ref{sec:dictionary}, the monomial-code classification.  These
claims are Proposition~\ref{prop:pencil-projective-recovery}, whose proof in
Appendix~\ref{sec:pencil-calculations} lists all four projectivities and
checks the bracket invariant.

For prime \(q\), Theorem~\ref{thm:minimum-support-transitions} assigns a
two-dimensional stabilizer-label space to every four-coordinate support.
Projection identifies that space with the Pauli-label plane at each retained
coordinate.  Composing one such identification with the inverse of another
gives a transition map, and closed chains of transitions give endomorphisms
\(M\) of a local Pauli-label plane.  Their projective conjugacy invariant is
\[
 r(M)=\Tr(M)^2/\det M,
\]
which is unchanged by conjugation, scalar multiplication, and inversion.
Its multiplicity data recover \(-1/z\), so local-Clifford equivalence also
forces equality of \(z\).  This is
Proposition~\ref{prop:pencil-holonomy-recovery}; its proof in
Appendix~\ref{sec:pencil-calculations} defines the transition maps, and
Appendix~\ref{sec:holonomy-certificate} records the exact collision
bookkeeping and certificate boundary.  Projective equivalence directly
supplies a local-Clifford equivalence.
\end{proof}

\begin{proof}[Proof of Corollary~\ref{cor:lu-lc-pencil}]
Theorem~\ref{thm:lc-pencil} gives all equivalences except local-unitary
equivalence.  Every local Clifford is local unitary, while
Theorem~\ref{thm:lu-lc-rigidity} turns every product-unitary intertwiner
between \(\Psi_t\) and \(\Psi_u\) into a product-Clifford intertwiner.
\end{proof}

The projective quotient remains valid over odd extension fields, but the
local-Clifford and local-unitary conclusion is intentionally restricted to
prime fields.  Frobenius maps are already additional local Clifford
operations, and \(z(t^p)=z(t)^p\) need not equal \(z(t)\).
Appendix~\ref{sec:pencil-calculations} records the Frobenius-sector divisors
and a concrete \(\F_{25}\) example; no complete extension-field orbit
classification is claimed.
